\par\noindent\SourceLeaf{59}%
oder kürzer
\[
= \frac{1}{2^{2h}}\,{}^{\nu}\mathfrak{B}^{(2h)} .  {}^{2h}\mathfrak{B}^{(h)}\,\frac{\nu+2}{h+1}\,n^{2h+1}.
\]
Setzen wir ferner $n$ als veränderlich an und bilden $\int An\,\mathrm{d}n$. Das Anfangsglied ist $=\frac{1}{2}nn$, das allgemeine Glied $h^{\text{tus}}$ ist
\[
= \frac{1}{2^{2h}}\,{}^{\nu}\mathfrak{B}^{(2h)} .  {}^{2h}\mathfrak{B}^{(h)}\,\frac{1}{2h+2}\,n^{2h+2}.
\]
Multipliziert man diese neue Reihe mit
\[
\frac{2(\nu+2)}{n}
\]
so liefert sie das Anfangsglied $=(\nu+2)n$ und das allgemeine Glied $h^{\text{tum}}$
\[
= \frac{1}{2^{2h}}\,{}^{\nu}\mathfrak{B}^{(2h)} .  {}^{2h}\mathfrak{B}^{(h)}\,\frac{\nu+2}{h+1}\,n^{2h+1}.
\]
Da die Glieder auf beiden Seiten übereinstimmen, ergibt sich
\[
2n .  A + A^{(1)} = \frac{2 . (\nu+2)}{n}\int An\,\mathrm{d}n,
\]
wobei das Integral ohne additive Konstante genommen wird. Daraus folgt
\[
A^{(1)} = \frac{2 . (\nu+2)}{n}\int An\,\mathrm{d}n - 2n .  A.
\]
\noindent\hspace*{\parindent}Mit dieser Formel wird also $A^{(1)}$ auf das bereits gefundene $A$ zurückgeführt, und nur dieses bleibt noch aus der Entstehung der Reihe selbst herzuleiten, wie es zuvor geschehen ist.

8. Nachdem erreicht ist, was bei der allgemeinsten Fassung der Aufgabe zu erwarten war, wenden wir uns den Sonderfällen zu. Denn in diesen lassen sich die allgemeinen Formeln bisweilen vereinfachen. Ein Beispiel ist der Fall, dass $\nu$ eine positive ganze Zahl ist. Dann muss nämlich jede der zuvor für die Koeffizienten von $(1+n\cos\varphi)^{\nu}$ gefundenen Reihen abbrechen, da jede in ihrer Reihenfolge die zum Exponenten $\nu$ gehörenden Binomialkoeffizienten enthält. Damit wird der Ausdruck für jeden Koeffizienten endlich. Ebenso kann man, wenn $\nu$ eine negative ganze Zahl ist, zu endlichen Ausdrücken für die einzelnen Koeffizienten gelangen. Da die Formel hauptsächlich in diesem Fall verwendet wird, lohnt es sich, bei seiner Ausarbeitung etwas zu verweilen. Das dazu nötige Verfahren gliedert sich in zwei Teile. Zuerst sind für $\nu=-1$ endliche Ausdrücke für die Koeffizienten zu finden. Dann ist zu zeigen, wie man bei gegebenen Koeffizienten von $(1+n\cos\varphi)^{\nu}$ die Koeffizienten der um einen Grad niedrigeren Potenz
